\VerdierPhysicalPageAnchor{0062}%
\VerdierSemanticAnchor{ch1-1-6-4}
\VerdierSemanticAnchor{ch1-1-6-4-definition}
\noindent\textbf{Definition 1.6.4.}\quad
The \emph{Koszul cocycle relative to an ordered basis
$e_1,\ldots,e_n$ of $\mathbb Z^n$} is the unique cocycle $\alpha$ in
the Koszul class relative to the basis $e_1,\ldots,e_n$ satisfying:

\VerdierSemanticAnchor{ch1-1-6-4-eq1}
\begin{equation}
  \alpha(\sigma_1,\sigma_2)=1
  \quad
  \forall\sigma_1\in\mathbb Z^n(\leq i),\quad
  \forall\sigma_2\in\mathbb Z^n(\geq i),\quad
  \forall i\in[1,n].
  \tag{1.6.4.1}
\end{equation}

\VerdierSemanticAnchor{ch1-1-6-4-coordinate-intro}
Let
$\sigma_1=\sum_{i=1}^{n}p_i e_i$ and
$\sigma_2=\sum_{i=1}^{n}q_i e_i$
be two elements of $\mathbb Z^n$; then:

\VerdierSemanticAnchor{ch1-1-6-4-eq2}
\begin{equation}
  \alpha(\sigma_1,\sigma_2)
  =(-1)^{\left(\sum_{i<j}q_i p_j\right)}.
  \tag{1.6.4.2}
\end{equation}

\VerdierSemanticAnchor{ch1-1-6-4-justification}
This definition requires justification. The existence of a cocycle in
the Koszul class satisfying (1.6.4.1) is clear: an immediate calculation
shows that the cocycle (1.6.4.2) satisfies (1.6.4.1) and belongs to the
Koszul class (1.6.2). To verify uniqueness, it suffices to observe that
every cocycle cohomologous to the cocycle $\alpha$ is of the form

\VerdierSemanticAnchor{ch1-1-6-4-cohomologous-form}
\begin{equation*}
  \alpha_1(\sigma_1,\sigma_2)
  =\alpha(\sigma_1,\sigma_2)
   \beta(\sigma_2)\beta(\sigma_1+\sigma_2)^{-1}\beta(\sigma_1)
\end{equation*}
and such a cocycle has property (1.6.4.1) if and only if

\VerdierSemanticAnchor{ch1-1-6-4-condition}
\begin{equation*}
  \beta(\sigma_1+\sigma_2)=\beta(\sigma_1)\beta(\sigma_2)
  \quad
  \forall i,\quad
  \forall\sigma_1\in\mathbb Z^n(\leq i),\quad
  \forall\sigma_2\in\mathbb Z^n(\geq i).
\end{equation*}

\VerdierSemanticAnchor{ch1-1-6-4-uniqueness}
It follows immediately that
$\beta\colon\mathbb Z^n\longrightarrow\mathbb Z/2\mathbb Z$
is a group homomorphism and consequently that
$\alpha_1=\alpha$.

\VerdierSemanticAnchor{ch1-1-6-5}
\VerdierSemanticAnchor{ch1-1-6-5-order-dependence}
\noindent\textbf{1.6.5.}\quad
The definition of the Koszul cocycle depends essentially on the total
order placed on the basis $e_1,\ldots,e_n$ of $\mathbb Z^n$. Let
$\leq$ denote the natural order on $[1,n]$, and let $\leq_s$ be a total
order on $[1,n]$. Let $s$ be the unique permutation of $[1,n]$ such
that:

\VerdierSemanticAnchor{ch1-1-6-5-eq1}
\begin{equation}
  i\leq_s j\quad\Longleftrightarrow\quad s(i)\leq s(j).
  \tag{1.6.5.1}
\end{equation}

\VerdierSemanticAnchor{ch1-1-6-5-cocycles}
Let $\alpha$ and $\alpha_s$ be the Koszul cocycles relative,
respectively, to the orders $\leq$ and $\leq_s$. Then:

\VerdierSemanticAnchor{ch1-1-6-5-eq2}
\VerdierSemanticAnchor{ch1-1-6-5-eq2-natural}
\VerdierSemanticAnchor{ch1-1-6-5-eq2-coordinates}
\VerdierSemanticAnchor{ch1-1-6-5-eq2-permuted}
\begin{equation}
  \begin{aligned}
    \alpha(\sigma_1,\sigma_2)
      &=(-1)^{\left(\sum_{i<j}q_i p_j\right)},\\
    \sigma_1&=\sum_i p_i e_i,
      \qquad \sigma_2=\sum_i q_i e_i,\\
    \alpha_s(\sigma_1,\sigma_2)
      &=(-1)^{\left(\sum_{s(i)<s(j)}q_i p_j\right)},
  \end{aligned}
  \tag{1.6.5.2}
\end{equation}
\VerdierSemanticAnchor{ch1-1-6-5-open}
